% ============================================================
% resume.tex — 由 AI 自动生成，基于 src/resume.md + templates/
% 请勿手动编辑此文件，修改简历内容请编辑 src/resume.md
% 修改样式请编辑 templates/style.sty
% ============================================================

\documentclass[10pt,a4paper]{article}

% --- 加载样式与模板 ---
\input{templates/style-aiagent.sty}
\input{templates/cv.latex}

% --- 局部扩展：四列布局等（调格式直接改数字） ---
\let\cvheadingThree\cvheading% 保存3列版本，项目部分恢复用
% 教育经历专用四列（时间 / 学校 / 学历 / 专业）
\newcommand{\cvheadingEdu}[4]{%
    \noindent
    \makebox[0.32\textwidth][l]{\textbf{#1}}% 时间
    \makebox[0.33\textwidth][l]{\textbf{#2}}% 学校
    \makebox[0.12\textwidth][l]{\textbf{#3}}% 学历
    \makebox[0.23\textwidth][r]{\textbf{#4}}\par% 专业
}
\newcommand{\cveventEdu}[5]{%
    \vspace{0.3em}
    \cvheadingEdu{#1}{#2}{#3}{#4}
    \ifx\relax#5\relax\else
        \vspace{0.06em}
        #5
    \fi
}
% 五列版本：工作经历用（时间 / 公司 / 领域 / 语言 / 岗位）
\newcommand{\cvheadingWork}[5]{%
    \noindent
    \makebox[0.28\textwidth][l]{\textbf{#1}}% 时间
    \makebox[0.18\textwidth][l]{\textbf{#2}}% 公司（图标+名称）
    \makebox[0.11\textwidth][l]{\textbf{#3}}% base地
    \makebox[0.23\textwidth][l]{\textbf{#4}}% 行业领域
    \makebox[0.20\textwidth][r]{\textbf{#5}}\par% 岗位+语言
}
% 工作经历专用 \cvevent（6参：时间/公司/base地/领域/岗位+语言/描述）
\newcommand{\cveventWork}[6]{%
    \vspace{0.2em}
    \cvheadingWork{#1}{#2}{#3}{#4}{#5}
    \ifx\relax#6\relax\else
        \vspace{0.04em}
        #6
    \fi
}
% section 标题旁的技术栈图标（PNG 彩色小图标）
% 技术栈图标：优先 devicon(-original) → simple-icons(-simple) → 手动(-custom)
\newcommand{\ticon}[1]{%
    \IfFileExists{assets/myIcons/#1-original.png}{%
        \raisebox{-0.08cm}{\includegraphics[height=0.42cm]{assets/myIcons/#1-original.png}}%
    }{\IfFileExists{assets/myIcons/#1-simple.png}{%
        \raisebox{-0.08cm}{\includegraphics[height=0.42cm]{assets/myIcons/#1-simple.png}}%
    }{%
        \raisebox{-0.08cm}{\includegraphics[height=0.42cm]{assets/myIcons/#1-custom.png}}%
    }}%
}
\newcommand{\seclogo}[1]{\hspace{0.12em}\ticon{#1}}

% ============================================================
% 个人信息设置
% ============================================================
\name{赵志文}
\lineA{%
    \faPhone\ 185-5613-5826 \enspace|\enspace
    \faEnvelope\ \href{mailto:cnjszzw@qq.com}{cnjszzw@qq.com} \enspace|\enspace
    \faWeixin\ cnjszzw%
}
\lineB{籍贯：江苏扬州 \enspace|\enspace 年龄：26 (1999/10/27) \enspace|\enspace 性别：男}
\lineC{求职意向：软件开发}
\photo{assets/photo.jpg}

% ============================================================
% 文档正文
% ============================================================
\begin{document}

\makecvheader

% ============================================================
% 教育经历
% ============================================================
\vspace{-0.825cm}%  第一个标题向上偏移，利用 header 下方空白
\cvsection{教育经历}

% 切换为4列布局（学校 / 学历 / 专业·学院·地点 / 时间）
\let\cvheading\cvheadingEdu
\let\cvevent\cveventEdu

\cvevent{2022/09 -- 2025/01}{%
    \raisebox{-0.12cm}{\includegraphics[height=0.50cm]{assets/school/Xijing_University_Emblem.png}}%
    \hspace{0.4em}西京学院}{硕士}{计算机技术}{}

\cvevent{2018/09 -- 2022/06}{%
    \raisebox{-0.12cm}{\includegraphics[height=0.50cm]{assets/school/Shandong_Normal_University_seal.png}}%
    \hspace{0.4em}山东师范大学}{本科}{广播电视编导}{}

% ============================================================
% 工作经历
% ============================================================
\cvsection{工作经历}

% 切换为4列布局（时间 / 公司 / 领域 / 岗位）
\let\cvheading\cvheadingWork
\let\cvevent\cveventWork

\cvevent{2025/04 -- 2026/03 (正职)}{%
    \raisebox{-0.12cm}{\includegraphics[height=0.55cm]{assets/LOGO/上海环创通信科技股份有限公司.png}}%
    \hspace{0.4em}寰创通信}{上海}{物联网(政企B端)}{(Java/Vue/Python)全栈开发}{}

\cvevent{2024/02 -- 2024/05 (实习)}{%
    \raisebox{-0.12cm}{\includegraphics[height=0.55cm]{assets/LOGO/美团.png}}%
    \hspace{0.4em}美团}{上海}{电商(B端商家侧)}{(Java)测试开发}{}

\cvevent{2023/09 -- 2024/02 (实习)}{%
    \raisebox{-0.12cm}{\includegraphics[height=0.55cm]{assets/LOGO/哔哩哔哩.png}}%
    \hspace{0.4em}哔哩哔哩}{上海}{内容社区(C端)}{(Python/Go)测试开发}{}

% ============================================================
% 项目经历（按公司分组）
% ============================================================
\let\cvheading\cvheadingThree% 恢复3列版本

\cvsection{上海寰创通信股份有限公司}[%
    \seclogo{claudecode}\seclogo{cursor}%                                  AI 工具
    \seclogo{python}\seclogo{langchain}\seclogo{langgraph}\seclogo{deepseek}\seclogo{mcp}\seclogo{fastapi}% Python/AI 栈
    \seclogo{java}\seclogo{spring}\seclogo{springboot}\seclogo{mybatis}\seclogo{mybatis-plus}\seclogo{mysql}\seclogo{redis}\seclogo{mqtt}% Java 栈
    \seclogo{vuejs}\seclogo{javascript}\seclogo{typescript}\seclogo{cesium}% 前端
    \seclogo{git}\seclogo{docker}\seclogo{jenkins}\seclogo{subversion}%    基础设施
]

\cvproject{项目一：无人机管理云平台}{}{2025/04 -- 2026/03}{%
    \cvtech{Java、SpringBoot、Redis、MySQL、MyBatis(Plus)、WebSocket、MQTT、Vue、JavaScript、TypeScript、Cesium}
    \cvlabeled{项目概述}{在公司物联网平台基础上，利用自组网超远距离（>20km）优势新增无人机接入模块，实现空地一体化指挥调度。}
    \cvlabeled{项目实现}{%
        \begin{itemize}
            \item \textbf{无人机接入与设备管理}
              \begin{starlist}
                \item[(S)背景:]无人机通过自组网硬件连接至基站，需作为新型设备纳入现有物联网平台统一管理。与地面设备不同，其连接链路更长、断连场景更复杂，现有上下线流程无法直接适配。
                \item[(T)任务:]将无人机融入现有平台的设备管理体系，复用权限管控与上下线流程，并解决服务重启后设备无法自动恢复在线的问题。
                \item[(A)动作:]后端将无人机接入现有设备管理流程，通过 Redis 管理在线状态，区分主动下线与异常断连。利用 Spring 启动就绪事件监听，服务启动后通过 MQTT 广播自定义重连指令，触发无人机重新走上线流程。
                \item[(R)结果:]无人机完全融入现有设备管理体系，与地面设备共享权限管控与上下线逻辑，服务重启后自动恢复上线，无需人工值守。
              \end{starlist}

            \item \textbf{飞行控制与相机系统}
              \begin{starlist}
                \item[(S)背景:]无人机飞行与相机指令类型繁多（涵盖一键起飞、指点飞行、环绕飞行、返航、急停等飞行指令，以及变焦、拍照、录像、云台重置、框选变焦、AIM 瞄准等相机指令）。前期为快速验证功能，大部分指令由前端通过 MQTT 直接与硬件通信，虽然开发快但指令逻辑分散在前端组件中，不利于维护和安全管控。
                \item[(T)任务:]完成全部飞行与相机指令的功能实现，将指令通道从前端直连逐步迁移至后端统一处理，并增加飞行安全校验防止误操作。
                \item[(A)动作:]实现一键起飞、指点飞行、环绕飞行、航线飞行、返航、急停等飞行指令，以及变焦、拍照、录像、云台重置、框选变焦、AIM 瞄准等相机负载指令。前期前端直连 MQTT 快速验证功能可行性，后期将指令逻辑迁移至后端统一处理。后端增加飞行安全校验：限高对比校验、降落确认与取消机制。
                \item[(R)结果:]覆盖全部飞行与相机操控能力，指令通道统一收敛至后端，飞行安全校验有效防止超限飞行等误操作风险。
              \end{starlist}

            \item \textbf{无人机直播观看}
              \begin{starlist}
                \item[(S)背景:]硬件推流不稳定，视频流可能延迟到达甚至不可达；多指挥员需同时观看同一无人机画面。
                \item[(T)任务:]在不可靠硬件条件下，实现可靠的多人直播观看，同时避免无人观看时浪费带宽。
                \item[(A)动作:]基于 Spring DeferredResult 实现异步确认机制——用户发起观看后立即释放线程，后端异步轮询 ZLM 等待推流就绪，确认可用后回调返回播放地址。同时设计 Redis 引用计数策略，维护观看用户集合，人数归零时自动停止拉流。
                \item[(R)结果:]消除了硬件延迟导致的线程阻塞，多人观看互不影响，无人观看时带宽自动回收。
              \end{starlist}
        \end{itemize}
    }
}

\cvproject{项目二：WVP Copilot --- 基于 LLM 的无人机 AI Agent 智能体}{}{2025/11 -- 2026/03}{%
    \cvtech{Python、LangChain、LangGraph、SSE、MCP、FastAPI、DeepSeek API、Redis}
    \cvlabeled{项目概述}{面向消防应急指挥场景，深度结合软硬件能力与 AI 算法打造无人机智能体 Copilot，极大降低使用门槛。用户只需自然语言描述需求（如"飞到陆家嘴，变焦 7x 拍照，返航"），Copilot 即可精准理解意图，自动拆分为无人机指令序列并编排执行，通过对话的方式完成远程控制与任务下达。采用 Python（Agent 编排层）+ Java（设备控制层）两层架构，通过 REST 通信解耦。}
    \cvlabeled{项目实现}{%
        \begin{itemize}
            \item \textbf{异步操作的状态同步与确定性故障恢复}
              \begin{starlist}
                \item[(S)背景:]LLM 推理为同步过程（秒级），无人机飞行为异步过程（分钟级）。工具函数调用飞行指令后需等待实际执行完成，但等待时间受飞行距离与通信延迟影响，不确定。
                \item[(T)任务:]在 LLM 同步推理与无人机异步执行之间建立可靠的状态同步，超时后需确定性判断飞行是否成功，避免误判导致后续指令出错。
                \item[(A)动作:]工具函数内部封装 Redis 轮询等待，依据飞行距离动态估算耗时阈值。超时后不直接判定失败，先查询 OSD 实际位置进行确定性兜底——距目标点 < 5m 判定已到达（通信延迟导致状态滞后）、位置无偏移判定硬件故障、位置持续变化则延长等待。中间状态仅向前端推送，LLM 不感知。
                \item[(R)结果:]异步等待的不确定性被消除，超时后通过 OSD 位置实现确定性判定，LLM 始终基于准确状态做出下一步决策。
              \end{starlist}

            \item \textbf{多步任务链的级联故障阻断与安全隔离}
              \begin{starlist}
                \item[(S)背景:]用户指令可能拆分为多步任务链（飞到某地 \thickarrow 拍照 \thickarrow 返航），链上任一步失败若不阻断，后续步骤可能在异常状态下执行造成安全风险。急停等安全操作需与 LLM 完全隔离，保证绝对可靠。
                \item[(T)任务:]实现多步任务链的故障阻断，任意步骤失败立即停止后续执行；高风险操作需人工确认；急停通道不经过 Agent 或 LLM。
                \item[(A)动作:]Plan-Confirm-Execute 三段式流程中，任意步骤返回失败或触发急停时立即 break 阻断后续执行。高风险操作（飞行、返航）执行前强制 Human-in-the-Loop 确认。急停指令走独立 MQTT 硬件通道，不经过 Agent 或 LLM，前端同步关闭 SSE 连接实现会话级隔离。
                \item[(R)结果:]多步任务链具备级联故障阻断能力，高风险操作有人工兜底，急停与 LLM 物理隔离保证绝对可靠。
              \end{starlist}

            \item \textbf{状态机驱动的子任务编排替代 LLM 时序决策}
              \begin{starlist}
                \item[(S)背景:]环绕飞行是天然状态机（进入环绕 \thickarrow 逐角度拍照 \thickarrow $360^\circ$ 退出），但 LLM 推理延迟（1-3s/次）在无人机 3m/s 飞行速度下将导致拍照点偏移 3-9 米。ReAct 循环依赖 LLM 判断时序，在实时控制场景下物理上不可行。
                \item[(T)任务:]将环绕飞行等实时控制任务的时序决策权从 LLM 移交给确定性状态机，同时保持 LLM 的意图理解能力。
                \item[(A)动作:]引入 LangGraph 建模显式条件子图——enter\_orbit \thickarrow orbit\_loop（条件边: $\ge 360^\circ$ 退出）\thickarrow exit\_orbit，执行周期由图边决定，LLM 仅负责意图路由而非时序判断。主链路（线性任务）保持 AgentExecutor 避免过度设计。
                \item[(R)结果:]消除 LLM 推理延迟导致的拍照点偏移，环绕时序精度由图状态转移保证，LLM 职责收窄为意图理解。
              \end{starlist}

            \item \textbf{MCP 外部工具链与复合指令封装}
              \begin{starlist}
                \item[(S)背景:]无人机控制涉及多个外部系统（高德地图地理编码、设备控制 API），单一用户意图（如"录像 30 秒"）往往映射为多个底层指令序列（start \thickarrow 倒计时 \thickarrow stop）。LLM 直接编排底层指令容易出错且低效。
                \item[(T)任务:]集成外部工具链，封装复合指令降低 LLM 编排复杂度，让 LLM 面向业务语义而非硬件原语编程。
                \item[(A)动作:]通过 MCP 协议集成高德地图 Server，15 个地图工具自动注册，地理编码多候选项回传 LLM 根据上下文仲裁最佳匹配。基于 LangChain @tool 定义 16 个无人机工具函数，遵循"暴露业务语义而非硬件原语"原则封装复合指令（如 record\_for\_duration 内部编排 start \thickarrow 倒计时 \thickarrow stop 完整时序）。SSE 流式推送实现工具调用链路的实时可观测。
                \item[(R)结果:]LLM 面向业务语义调用工具，复合指令将常见操作从多步简化为单步，SSE 流式输出让操作过程对用户透明。
              \end{starlist}
        \end{itemize}
    }
}

\cvsection{个人开源项目}[\seclogo{claudecode}\seclogo{java}\seclogo{spring}\seclogo{springboot}\seclogo{mybatis}\seclogo{mysql}\seclogo{redis}\seclogo{git}\seclogo{elasticsearch}\seclogo{rocketmq}\seclogo{nginx}]

\cvproject{项目一：仿"B站"弹幕视频网站}{}{2024/05 -- 2026/05}{%
    \cvlabeled{项目地址}{%
        \begingroup
        \hypersetup{urlcolor=primaryblue}%
        \url{https://github.com/Cnjszzw/imooc-bilibili} \textbar\ \url{https://github.com/Cnjszzw/bilibili-microservices}%
        \endgroup
    }
    \cvtech{SpringCloud Alibaba、Nacos、Gateway、OpenFeign、Seata、MyBatis、Redis、RocketMQ、Elasticsearch、WebSocket}
    \cvlabeled{项目概述}{实现B站核心功能——用户体系、视频管理、弹幕系统、站内搜索，涵盖从注册登录、视频上传观看、弹幕实时推送到全文检索的完整闭环。}
    \cvlabeled{项目实现}{%
        \begin{itemize}
            \item \textbf{用户体系与安全}：基于工厂 + 策略模式构建登录系统（UserGranter 策略接口 + UserLoginFactory 配置驱动注册），新增登录方式只需实现接口并配置，无需修改调用方代码。基于 JWT 双 Token（access + refresh）实现无感登录，RSA 加密传输 + MD5 加盐存储保证密码安全。基于 RBAC + AOP 切面实现接口级和数据级权限控制，与业务代码解耦。

            \item \textbf{站内全文搜索}：集成 Elasticsearch，基于 Multi-Match 实现跨字段检索（title\slash nick\slash description），支持关键词高亮、多维度排序（时间/弹幕数/播放量）及分页查询。用户注册和视频发布时同步写入 ES 索引，保证搜索实时性。

            \item \textbf{视频存储与传输}：集成 FastDFS 实现视频分片上传、断点续传与秒传（文件指纹校验），通过 Nginx 提供外部 HTTP 访问。

            \item \textbf{弹幕与动态推送}：基于 WebSocket 构建弹幕实时推送，弹幕消息先写入 RocketMQ 异步消费落库，避免高并发下直接写 DB 导致压力过大。动态发布系统经历三轮迭代：V1 同步推 \ \thickarrow\  V2 MQ 异步推 \ \thickarrow\  V3 推拉结合。推拉结合思路：粉丝 < 10w 走推模式（遍历写入粉丝收件箱），粉丝 $\ge$ 10w 的头部 UP 主走拉模式（只写自己发件箱，粉丝读时主动拉取），消除大 V 写扩散问题，支撑百万级粉丝场景。经压测验证，10w 粉丝场景下遍历写入从 7s 降至 $\sim$1ms，用户发布等待始终维持在 40ms 以内。

            \item \textbf{微服务架构改造}：基于 Spring Cloud Alibaba 完成单体应用微服务化升级，搭建 Nacos 注册中心 + Spring Cloud Gateway 统一网关，通过全局过滤器实现 JWT 统一鉴权与 X-User-Id 请求头注入，下游服务无需各自处理鉴权逻辑；采用 OpenFeign + Nacos 服务发现实现跨服务调用。

            \item \textbf{分布式事务}：视频投稿涉及 content-service 写视频表和 legacy-service 创建动态两个跨服务操作。评估 XA / Seata AT / TCC / 最大努力通知四种方案后，选择 RocketMQ 事务消息实现最终一致性：content-service 本地事务内写视频 + 发半消息，Broker 回查确认后投递，legacy 异步消费创建动态。设计入库消费者与扩散消费者独立重试（16 次递增），解决 Seata AT 无法覆盖 MQ + Redis 的跨存储回滚问题，实现 DB / MQ / Redis 三层最终一致。
        \end{itemize}
    }
}

% ============================================================
% 专业技能
% ============================================================
\cvsection{专业技能}

\begin{multicols}{2}
\cvskill{Java 全家桶}{熟练掌握 Java，熟悉 JVM 内存模型与并发编程；熟练使用 SpringBoot、MyBatis（Plus）；了解 SpringCloud 微服务体系（Nacos、OpenFeign、Gateway）}

\cvskill{数据存储与检索}{熟悉 MySQL，有 N+1 查询优化（预加载 + HashMap 索引）、延迟关联、覆盖索引等 SQL 调优实战经验；熟悉 Redis 缓存策略与数据一致性保障；熟悉 Elasticsearch 全文检索（Multi-Match、高亮、多维度排序）}

\cvskill{AI Agent 工程化}{有 LLM Agent 从 0 到 1 的实践经验：熟练使用 LangChain、LangGraph、Function Calling、MCP 协议；掌握 Agent 可靠性设计；熟练使用 Cursor/ClaudeCode}

\cvskill{分布式系统}{了解分布式锁（Redisson）、分布式事务（RocketMQ 事务消息 / TCC / 可靠消息最终一致）、配置中心（Nacos）等常见组件的设计思路与应用场景}

\cvskill{性能优化}{有核心接口性能优化实战经验：千台 Mesh 设备查询从 7-8s 优化至 1s 以内，30w 条日志深分页从 5s 降至 1s 以内（性能提升 5 倍）；善于定位并解决大数据量场景下的查询与存储瓶颈}

\cvskill{工程实践}{熟悉 Logback 日志框架与 Kibana 排障，熟悉 Linux + Docker 基础运维；具备从需求分析到联调交付的全流程能力，有前后端 + 软硬件跨团队协作经验}
\end{multicols}

% ============================================================
% 竞赛/获奖/证书
% ============================================================
\cvsection{竞赛/获奖/证书}

\begin{itemize}
    \item 考研初试368分（专业第一）\enspace|\enspace 2022西京学院一等奖学金（前20\%）\enspace|\enspace 2022传智杯程序设计赛区域赛二等奖（前20\%）
    \item 雅思IELTS（6.5分）\enspace|\enspace 英语六级CET-6（593分）\enspace|\enspace 国际人才英语考试ETIC（优秀等级）\enspace|\enspace 计算机二级证书（MS OFFICE）
\end{itemize}

\end{document}
